% Complete uniform-integrability and conditional-limit proofs.

\begin{annexitem}{annex:p15.ui-small-sets}{thm:p15-ui-small-sets}{Tail splitting and small events}
\begin{proof}
An empty family satisfies every assertion, so suppose the family is nonempty.
Uniform integrability gives an \(M>0\) with
\(\sup_iE(|X_i|;|X_i|>M)\leq1\). Hence \(E|X_i|\leq M+1\).
For any event \(A\) and any \(M>0\),
\[
             E(|X_i|;A)\leq MP(A)+E(|X_i|;|X_i|>M).
\]
Given \(\epsilon>0\), choose \(M>0\) making the last term at most
\(\epsilon/2\) for every \(i\), and take
\(\delta=\epsilon/(2M)\). This proves uniform small-event control.

Conversely, let \(C=\sup_iE|X_i|<\infty\). If \(C=0\), all variables
are zero almost surely. If \(C>0\), Markov's inequality gives
\(P(|X_i|>M)\leq C/M\). For the \(\delta\) provided by small-event
control, take \(M\geq C/\delta\) and apply that control to
\(\{|X_i|>M\}\). This proves uniform integrability.

If \(|X_i|\leq Z\) almost surely for an integrable \(Z\geq0\), then
\[
      E(|X_i|;|X_i|>M)\leq E(Z;Z>M)\longrightarrow0
\]
by dominated convergence. For a finite family one may use
\(Z=\sum_i|X_i|\), which is integrable. Finally
\(X_n=n\mathbf1_{(0,1/n)}\) on the uniform unit interval has mean one
and, for every \(M\), tail mean one whenever \(n>M\). It establishes
the failure of the first-moment condition alone.
\end{proof}
\end{annexitem}

\begin{annexitem}{annex:p15.ui-moment-criteria}{thm:p15-ui-moment-criteria}{The source bound and an explicit convex converse}
\begin{proof}
For each \(\epsilon>0\), superlinear growth gives \(M_\epsilon\) with
\(t\leq\epsilon\Phi(t)\) for all \(t>M_\epsilon\).
Thus, writing \(C=\sup_iE\Phi(|X_i|)<\infty\),
\[
   E(|X_i|;|X_i|>M_\epsilon)
       \leq\epsilon E(\Phi(|X_i|);|X_i|>M_\epsilon)
       \leq\epsilon C.
\]
This proves the source criterion; if \(C=0\) the displayed bound is
already zero. The two examples are nonnegative Borel functions and
their ratios to \(t\) tend to infinity.

For the converse choose finite numbers \(a_j\), \(j\geq1\), recursively
so that \(a_j\geq j\), \(a_{j+1}>a_j\), and
\[
                   \sup_iE(|X_i|;|X_i|>a_j)\leq2^{-j}.
\]
Uniform integrability permits every choice. Define
\(\Phi(t)=\sum_{j\geq1}(t-a_j)^+\) for \(t\geq0\).
On a bounded interval only finitely many terms are nonzero, so this
is finite and continuous. Each summand is convex and nondecreasing;
the same properties follow for their pointwise sum by summing the
defining inequalities. Also \(\Phi(0)=0\). For every integer \(r\),
\[
             \liminf_{t\to\infty}\frac{\Phi(t)}t
                  \geq\lim_{t\to\infty}
                       \sum_{j=1}^r\frac{(t-a_j)^+}t=r.
\]
Since \(r\) is arbitrary, the ratio tends to infinity.
Monotone convergence for the nonnegative partial sums yields
\[
 E\Phi(|X_i|)
     =\sum_{j\geq1}E(|X_i|-a_j)^+
     \leq\sum_{j\geq1}E(|X_i|;|X_i|>a_j)
     \leq\sum_{j\geq1}2^{-j}=1.
\]
For the empty family the explicit choice \(\Phi(t)=t^2\) suffices.
\end{proof}
\end{annexitem}

\begin{annexitem}{annex:p15.ui-conditional-family}{thm:p15-ui-conditional-family}{One absolute-continuity bound for every information set}
\begin{proof}
The singleton \(\{|X|\}\) has small-event control by the preceding
tail-splitting proof. Fix \(\epsilon>0\), and choose its \(\delta>0\).
Let \(V_{\mathcal G}=E(|X|\mid\mathcal G)\). This is nonnegative,
integrable, and has mean \(E|X|\). For \(M>0\) with \(E|X|/M\leq\delta\),
Markov's inequality gives \(P(V_{\mathcal G}>M)\leq\delta\).
Conditional Jensen gives \(|E(X\mid\mathcal G)|\leq V_{\mathcal G}\),
so, using the defining identity on the \(\mathcal G\)-event
\(\{V_{\mathcal G}>M\}\),
\[
 \begin{aligned}
 E\bigl(|E(X\mid\mathcal G)|;\,|E(X\mid\mathcal G)|>M\bigr)
  &\leq E(V_{\mathcal G};V_{\mathcal G}>M)\\
  &=E(|X|;V_{\mathcal G}>M)\leq\epsilon .
 \end{aligned}
\]
The same \(M\) works for every \(\mathcal G\).
When \(E|X|=0\), every conditional expectation is zero almost surely.
This proof uses the small-event integral estimate directly; convergence
of event probabilities does not require pointwise convergence of their
indicators.
\end{proof}
\end{annexitem}

\begin{annexitem}{annex:p15.probability-l1-tools}{thm:p15-probability-l1-tools}{All five elementary limit steps}
\begin{proof}
Markov's inequality gives
\(P(|X_n-X|>\epsilon)\leq E|X_n-X|/\epsilon\), proving the first assertion.
For the second choose strictly increasing indices \(n_k\) with
\[
                 P(|X_{n_k}-X|>2^{-k})\leq2^{-k}.
\]
The first Borel--Cantelli lemma makes all but finitely many of these
deviations absent almost surely, so \(X_{n_k}\to X\) there.

For the bounded assertion,
\[
                E|Z_n-Z|\leq\epsilon+2K P(|Z_n-Z|>\epsilon).
\]
Take the limit superior and then let \(\epsilon\downarrow0\).
If \(K=0\), the assertion is immediate.

For the fourth assertion there is nothing to prove if
\(\liminf E|X_n|=\infty\). Otherwise select a subsequence whose expectations
tend to that finite lower limit and then an almost-surely converging
further subsequence by the second assertion. Fatou's lemma applied to
its absolute values proves the inequality. The same argument covers a
zero lower limit.

Finally, for every measurable \(A\), absolute integrability and the
triangle inequality give
\[
 |E(X_n;A)-E(X;A)|
        \leq E(|X_n-X|\mathbf1_A)\leq E|X_n-X|.
\]
The right side is independent of \(A\), proving uniformity.
\end{proof}
\end{annexitem}

\begin{annexitem}{annex:p15.vitali-convergence}{thm:p15-vitali-convergence}{Both complete source truncations and every implication}
\begin{proof}
Assume uniform integrability. Its first-moment bound and
\cref{thm:p15-probability-l1-tools} imply \(E|X|<\infty\).
Define the clipping map
\(\phi_M(x)=\max(-M,\min(x,M))\), \(M>0\).
This map is 1-Lipschitz, so
\(\phi_M(X_n)\to\phi_M(X)\) in probability. Both clipped variables
are bounded by \(M\); their difference therefore tends to zero in
\(L^1\). Pointwise,
\[
 \begin{split}
 |X_n-X|\leq&
       |\phi_M(X_n)-\phi_M(X)|\\
       &+|X_n|\mathbf1_{\{|X_n|>M\}}
        +|X|\mathbf1_{\{|X|>M\}} .
 \end{split}
\]
For any \(\epsilon>0\), take \(M\) sufficiently large that both tail
expectations are at most \(\epsilon\), uniformly in \(n\) for the first.
For this fixed \(M\) the clipped expectation tends to zero. Thus
\(\limsup E|X_n-X|\leq2\epsilon\), proving convergence in \(L^1\).

The implication from \(L^1\) convergence to convergence of absolute
means follows from
\[
                 |E|X_n|-E|X||\leq E\bigl||X_n|-|X|\bigr|
                                      \leq E|X_n-X|.
\]

For the remaining implication assume \(E|X_n|\to E|X|<\infty\).
For \(M>1\) define the source cutoff, on \(t\geq0\), by
\[
 \psi_M(t)=
 \begin{cases}
 t,&0\leq t\leq M-1,\\
 (M-1)(M-t),&M-1<t<M,\\
 0,&t\geq M.
 \end{cases}
\]
It is continuous, bounded, piecewise linear, and \(0\leq\psi_M(t)\leq t\).
For each fixed \(M\) it is Lipschitz, so
\(\psi_M(|X_n|)\to\psi_M(|X|)\) in probability and in \(L^1\).
Also \(\psi_M(|X|)\to|X|\) as \(M\to\infty\), dominated by \(|X|\).
Choose \(M\) so that
\(E|X|-E\psi_M(|X|)<\epsilon/2\).
Convergence of the two expectations gives an \(n_0\) such that, for
\(n\geq n_0\),
\[
 \begin{aligned}
 E(|X_n|;|X_n|>M)
   &\leq E|X_n|-E\psi_M(|X_n|)\\
   &<E|X|-E\psi_M(|X|)+\epsilon/2<\epsilon.
 \end{aligned}
\]
The finitely many earlier variables are integrable. Increase the tail
threshold to \(M'\geq M\) until their tail expectations are also at most
\(\epsilon\). For \(n\geq n_0\) increasing the threshold preserves the
previous bound, regardless of how the cutoff function changes.
This establishes uniform integrability for the entire sequence.
\end{proof}
\end{annexitem}

\begin{annexitem}{annex:p15.submartingale-ui}{thm:p15-submartingale-ui}{Mean convergence and the terminal conditional inequality}
\begin{proof}
Uniform integrability bounds \(\sup_nE|X_n|\), hence the positive-part
moments. The proved martingale convergence theorem
\cref{thm:p15-martingale-convergence} gives an integrable almost-sure
limit. The event that the sequence has a finite limit is
\(\mathcal F_\infty\)-measurable: limsup and liminf of the measurable
variables are measurable. Define \(X_\infty\) to be the limit on that
event and zero elsewhere. Vitali's theorem gives \(L^1\) convergence.
Convergence in both senses implies convergence in \(L^1\). Conversely
\(L^1\) convergence gives convergence in probability, and Vitali's
theorem gives uniform integrability. This proves the equivalences.

For fixed \(n\), repeated conditional expectation yields
\(E(X_m\mid\mathcal F_n)\geq X_n\) for every \(m\geq n\).
Consequently \(\int_A X_m\,dP\geq\int_A X_n\,dP\) for \(A\in\mathcal F_n\).
Pass to the \(L^1\) limit using the preceding event-integral bound:
\(\int_A X_\infty\,dP\geq\int_A X_n\,dP\).
By the defining conditional identity this is
\(\int_A[E(X_\infty\mid\mathcal F_n)-X_n]\,dP\geq0\).
If the bracket were negative on a set of positive probability, one
of its \(\mathcal F_n\)-events where it is at most \(-1/k\)
would contradict that inequality. This proves the conditional inequality.
Apply the result to \(-X_n\) for a supermartingale.
\end{proof}
\end{annexitem}

\begin{annexitem}{annex:p15.closed-martingale}{thm:p15-closed-martingale}{Closing the martingale and identifying its observable terminal part}
\begin{proof}
The first three conditions are equivalent by the submartingale result.
If \(X_n\to X_\infty\) in \(L^1\), then for \(m\geq n\) and
\(A\in\mathcal F_n\) the martingale property gives
\(\int_A X_m\,dP=\int_A X_n\,dP\).
Taking the limit proves
\(\int_A X_\infty\,dP=\int_A X_n\,dP\).
Since \(X_n\) is integrable and \(\mathcal F_n\)-measurable, it is a
version of \(E(X_\infty\mid\mathcal F_n)\). Thus the limit closes it.
Conversely, the conditional-family theorem makes any martingale
closed by an integrable \(Z\) uniformly integrable.

To identify its limit for a given closing variable \(Z\), take the
\(\mathcal F_\infty\)-measurable limit version constructed above.
For every \(A\in\bigcup_n\mathcal F_n\), its integral equals that of \(Z\).
Let
\[
       \mathcal L=\{A\in\mathcal F_\infty:
                              \int_A X_\infty\,dP=\int_A Z\,dP\}.
\]
This is a lambda-system: it contains \(\Omega\), equality for a complement
follows by subtraction from the total integrals, and equality for a
countable disjoint union follows from dominated convergence of its
finite-union indicators, dominated separately by \(|X_\infty|\) and \(|Z|\).
The increasing union \(\bigcup_n\mathcal F_n\) is a pi-system,
since finitely many events belong to one sufficiently large sigma-algebra.
The proved pi-lambda theorem implies
\(\mathcal L=\mathcal F_\infty\). This is exactly the defining identity
for \(E(Z\mid\mathcal F_\infty)\), with the required measurability and
integrability. No unobserved part of \(Z\) is identified with the limit.
\end{proof}
\end{annexitem}

\begin{annexitem}{annex:p15.levy-upward}{thm:p15-levy-upward}{The full increasing-information argument}
\begin{proof}
Set \(M_n=E(Z\mid\mathcal F_n)\), choosing a finite real measurable
version at each time. Integrability and adaptedness follow from
conditional existence. The tower identity gives
\[
                    E(M_{n+1}\mid\mathcal F_n)=M_n.
\]
The conditional-family theorem makes \((M_n)\) uniformly integrable.
The closed-martingale theorem therefore gives an almost-sure and
\(L^1\) limit \(M_\infty\).

For completeness, the source's identification step is as follows.
The event of finite convergence is \(\mathcal F_\infty\)-measurable;
define \(M_\infty=0\) off that event. For any \(A\in\mathcal F_n\)
and \(m\geq n\),
\(\int_A M_m\,dP=\int_A Z\,dP\). The \(L^1\) limit preserves this equality.
The collection of all \(A\in\mathcal F_\infty\) on which the equality
holds is a lambda-system, by total integrals, complements and dominated
convergence on disjoint unions. It contains the pi-system
\(\bigcup_n\mathcal F_n\), and therefore contains its generated
sigma-algebra. The resulting identity and measurability identify
\(M_\infty=E(Z\mid\mathcal F_\infty)\) almost surely.
Neither limit construction requires complete information sigma-algebras.
\end{proof}
\end{annexitem}

\begin{annexitem}{annex:p15.levy-information-law}{thm:p15-levy-information-law}{Event learning and the separate tail-law proof}
\begin{proof}
The indicator of \(A\in\mathcal F_\infty\) is integrable and known at
that information level, so \(E(\mathbf1_A\mid\mathcal F_\infty)=\mathbf1_A\).
The upward theorem proves both claimed modes of convergence.

For the second assertion put
\(\mathcal F_n=\sigma(\xi_1,\ldots,\xi_n)\).
An event \(A\) in the tail field belongs to
\(\sigma(\xi_{n+1},\xi_{n+2},\ldots)\) for every \(n\), which is
independent of \(\mathcal F_n\) by generated independence.
Thus \(E(\mathbf1_A\mid\mathcal F_n)=P(A)\) almost surely for every \(n\).
Also \(A\in\sigma(\xi_1,\xi_2,\ldots)=\mathcal F_\infty\).
The countably many conditional identities hold together outside one
null set, and the first conclusion gives \(P(A)=\mathbf1_A\) there.
Taking the variance, or just observing the two possible values of an
indicator, gives \(P(A)(1-P(A))=0\).
\end{proof}
\end{annexitem}

\begin{annexitem}{annex:p15.changing-conditional-limits}{thm:p15-changing-conditional-limits}{Contraction and the full oscillation-envelope proof}
\begin{proof}
For the first assertion conditional \(L^1\) contraction gives
\[
 \begin{split}
 &E|E(Y_n\mid\mathcal F_n)-E(Y\mid\mathcal F_\infty)|\\
 &\quad\leq E|Y_n-Y|
       +E|E(Y\mid\mathcal F_n)-E(Y\mid\mathcal F_\infty)|
                      \longrightarrow0
 \end{split}
\]
by the upward theorem. Under domination, ordinary dominated convergence
already gives \(Y_n\to Y\) in \(L^1\), proving this part of the second assertion.

For its almost-sure assertion the countable supremum \(W_N\) is
measurable, is at most \(2Z\) almost surely, and decreases to zero
almost surely since a convergent real sequence is Cauchy.
For \(n\geq N\), taking the limit in \(m\) gives
\(|Y_n-Y|\leq W_N\). Conditional order implies
\[
                  E(|Y_n-Y|\mid\mathcal F_n)
                           \leq E(W_N\mid\mathcal F_n).
\]
For each fixed \(N\), the upward theorem identifies the limit of the
right side as \(E(W_N\mid\mathcal F_\infty)\).
There are countably many \(N,n\), so all these comparisons and limits
hold on one event of probability one. On that event,
\[
       \limsup_n E(|Y_n-Y|\mid\mathcal F_n)
                    \leq E(W_N\mid\mathcal F_\infty)
                                    \quad\text{for every }N.
\]
Conditional dominated convergence, with integrable bound \(2Z\),
makes the right side decrease to zero. Conditional Jensen and linearity
give
\[
 |E(Y_n\mid\mathcal F_n)-E(Y\mid\mathcal F_n)|
                         \leq E(|Y_n-Y|\mid\mathcal F_n)\longrightarrow0.
\]
Apply the upward theorem once more to \(Y\); the triangle inequality
proves the desired almost-sure limit.
\end{proof}
\end{annexitem}

\begin{annexitem}{annex:p15.ui-counterexamples}{thm:p15-ui-counterexamples}{Every law, limit and conditional identity in the examples}
\begin{proof}
In the first example, for every fixed \(x>0\), eventually \(1/n<x\),
so \(X_n(x)=0\). The interval has mass \(1/n\), hence \(E|X_n|=1\).
For every finite \(M\), choose \(n>M\); then
\(E(|X_n|;|X_n|>M)=1\), disproving uniform integrability.

In the second example all absolute values equal one, so the family
is dominated by the integrable constant one. Successive values differ
in absolute value by two with probability one. If the sequence
converged in probability to \(X\), the triangle inequality would imply
\[
 P(|X_{n+1}-X_n|>1)
 \leq P(|X_{n+1}-X|>1/2)+P(|X_n-X|>1/2)\longrightarrow0,
\]
a contradiction.

Construct the Bernoulli families in the third example on their countable
product space. At \(n=1\) both variables equal one; for every \(n\),
independence gives \(P(X_n>0)=1/n^2\) and \(EX_n=1/n\).
The first Borel--Cantelli lemma gives \(X_n\to0\) almost surely, and
the means give \(L^1\) convergence. For \(M>0\),
\[
 E(X_n;X_n>M)=
       \begin{cases}1/n,&n>M,\\0,&n\leq M.\end{cases}
\]
The supremum tends to zero, directly proving uniform integrability.
Generated independence makes \(C_n\) independent of the whole
\(\mathcal G=\sigma(B_1,B_2,\ldots)\). Therefore
\[
 E(nB_nC_n\mid\mathcal G)=nB_nEC_n=B_n.
\]
The \(B_n\) are independent and their success probabilities have divergent
sum, so the second Borel--Cantelli lemma gives infinitely many ones
almost surely. The zero events are also independent and their
probabilities \(1-1/n\), \(n\geq2\), have divergent sum, so there are
infinitely many zeros. Hence \(B_n\) has no almost-sure limit,
although \(E|B_n|=1/n\to0\).
\end{proof}
\end{annexitem}

\begin{annexitem}{annex:p15.dyadic-martingales}{thm:p15-dyadic-martingales}{Generating the Borel sets and passing telescoping increments to the limit}
\begin{proof}
Each dyadic atom is Borel and is the disjoint union of its two child
atoms. Thus \(\mathcal D_n\) increases. The family of all dyadic atoms
generates the Borel sets of \([0,1)\): for an open set \(O\) relative
to this interval and \(x\in O\), there is an \(\epsilon>0\) with
\((x-\epsilon,x+\epsilon)\cap[0,1)\subseteq O\).
At a sufficiently fine level the unique atom containing \(x\)
is contained in that neighborhood. Hence \(O\) is the union of the
dyadic atoms contained in it. There are only countably many atoms
over all levels, so this is a measurable union.

The finite-partition conditional formula gives the stated cell average,
since each atom has mass \(2^{-n}\). Apply the upward theorem to \(f\).
The limiting sigma-algebra is the Borel sigma-algebra, so its conditional
expectation of \(f\) equals \(f\). Each average is a finite dyadic step
function, and convergence in \(L^1\) supplies any prescribed approximation
error.

Here is the extension to completed Lebesgue measurability.
Choose finite-valued measurable functions \(s_m\) converging pointwise
to \(f\) wherever \(f\) is finite, for example by clipping to \([-m,m]\)
and rounding down to multiples of \(2^{-m}\).
Each of the finitely many level sets of each \(s_m\) differs from a
Borel set by a subset of a Borel null set, by the definition of
completion. Replacing its level sets by those Borel representatives
produces a Borel simple function \(t_m\) equal to \(s_m\) outside a
null set. The possible overlaps between representatives affect only
that null set. The countable union of these exceptional sets is null.
Define \(g\) as the finite limit of \(t_m\) where it exists and as zero
elsewhere. This is Borel measurable and equals \(f\) almost everywhere.
Its integrability, all dyadic cell integrals, and every \(L^1\)
approximation assertion are therefore the same as for \(f\).

For the Lipschitz assertion, any sequence \(t_j\uparrow1\) makes
\(f(t_j)\) Cauchy by the Lipschitz bound. Interlacing two such sequences
shows the limit is independent of the sequence; comparing arbitrary
\(t\to1\) with one chosen sequence gives the same limit.
Defining \(f(1)\) by it extends the original Lipschitz inequality to
all of \([0,1]\).
The slopes \(X_n\) are \(\mathcal D_n\)-measurable and satisfy
\(|X_n|\leq K\). If an atom at level \(n\) has endpoints \(l,r\)
and midpoint \(m\), its two child slopes have conditional average
\[
 \tfrac12\,2^{n+1}[f(m)-f(l)]
 +\tfrac12\,2^{n+1}[f(r)-f(m)]
                   =2^n[f(r)-f(l)].
\]
This is its parent slope, so the sequence is a martingale.
Domination by \(K\) implies uniform integrability, and its real limit
has \(|X_\infty|\leq K\) almost everywhere and convergence in \(L^1\).
Choose a bounded Borel version of that limit.

For dyadic endpoints \(a\leq b\), and all sufficiently large \(n\),
\([a,b)\) is a union of level-\(n\) atoms. Integrating \(X_n\) on them
telescopes to \(f(b)-f(a)\). Convergence in \(L^1\) therefore gives
\(\int_a^b X_\infty=f(b)-f(a)\) for these endpoints, including \(b=1\).
For arbitrary \(0\leq a<b\leq1\), take dyadic \(a_m\to a\),
\(b_m\to b\) with \(a_m\leq b_m\).
The change in the integral is at most
\(K(|a_m-a|+|b_m-b|)\), and continuity of \(f\) gives the limit of
the endpoint difference. For \(a=b\), both sides are zero.
This proves the identity for every stated pair of endpoints.
\end{proof}
\end{annexitem}

\begin{annexitem}{annex:p15.posterior-mean-consistency}{thm:p15-posterior-mean-consistency}{Identifying the parameter from the infinite observation}
\begin{proof}
The independent integrable strong law gives
\[
             \frac1n\sum_{i=1}^n\xi_i\longrightarrow\mu
                                      \quad\hbox{almost surely}.
\]
It follows on that same event that
\(n^{-1}\sum_{i=1}^nY_i\to\theta+\mu\).
Each observed average is \(\mathcal F_\infty\)-measurable.
Define \(L\) to be its finite limit where it exists, and zero elsewhere.
The convergence event and \(L\) are \(\mathcal F_\infty\)-measurable.
Thus \(L-\mu\) is an \(\mathcal F_\infty\)-measurable version of \(\theta\),
and \(E(\theta\mid\mathcal F_\infty)=\theta\) almost surely.
The upward theorem, using \(E|\theta|<\infty\), gives
the asserted almost-sure and \(L^1\) convergence.
The argument preserves the stated independent-prior sampling assumptions;
it does not require a particular prior density or Gaussian conjugacy.
For normal \((0,1)\) noise the needed integrability and mean follow
from the normal-law moment calculation.
\end{proof}
\end{annexitem}

\begin{annexitem}{annex:p15.extinction-or-escape}{thm:p15-extinction-or-escape}{The full event-learning argument and branching construction}
\begin{proof}
The event \(D\) belongs to \(\mathcal F_\infty\). L\'evy's event theorem
gives \(P(D\mid\mathcal F_n)\to\mathbf1_D\) almost surely.
Take one event of probability one on which this limit and all of the
countably many stated inequalities, over \(K,n\), hold.
At an outcome in this event but outside \(D\), the conditional
probabilities tend to zero. For fixed \(K\), infinitely many occurrences
of \(X_n\leq K\) would force a subsequence of those probabilities to
be at least \(\delta_K>0\), a contradiction. Thus eventually \(X_n>K\)
for every positive integer \(K\). This is exactly \(X_n\to\infty\).

For the branching application let the generation information be
\(\mathcal H_n=\sigma(Z_0,\{\xi_{r,j}:r<n,\ j\geq1\})\);
then \(Z_n\) is \(\mathcal H_n\)-measurable.
Each generation is finite almost surely by induction, since it is a
finite sum of finite integer offspring counts. The next row of offspring
variables is independent of \(\mathcal H_n\). On \(\{Z_n=m\}\),
all \(m\) offspring counts vanish with probability \(p_0^m\).
Conditioning and summing over the countably many possible \(m\) gives
\[
                       P(Z_{n+1}=0\mid\mathcal H_n)=p_0^{Z_n}.
\]
This identity includes \(m=0\), with the empty product equal to one.
Condition once more on \(\sigma(Z_0,\ldots,Z_n)\) to obtain the same
identity for the natural filtration.
On \(\{Z_n\leq K\}\), the chance of eventual zero is at least
\(p_0^{Z_n}\geq p_0^K>0\). Apply the first assertion to the shifted
sequence \(X_{n+1}=Z_n\), so that initial extinction is included.
Once a generation is zero, every later empty offspring sum is zero.
Hence hitting zero is equivalent to eventual zero, and the asserted
dichotomy follows. No first moment of \(Z_0\) or of an offspring count
was used.
\end{proof}
\end{annexitem}

\begin{annexitem}{annex:p15.bounded-polarization}{thm:p15-bounded-polarization}{Conditional means and the two possible terminal values}
\begin{proof}
The two stated possible next values differ by the strictly positive
constant \(\alpha\). Their conditional probabilities sum to one,
so they exhaust the conditional law. Therefore
\[
 E(X_{n+1}\mid\mathcal F_n)
      =X_n(\alpha+\beta X_n)+(1-X_n)\beta X_n
      =(\alpha+\beta)X_n=X_n.
\]
Thus the bounded sequence is a uniformly integrable martingale.
It has an almost-sure and \(L^1\) limit \(X_\infty\in[0,1]\),
and \(EX_\infty=EX_0\).
At every time \(X_{n+1}-\beta X_n\) belongs to \(\{0,\alpha\}\)
almost surely. All these memberships and convergence hold together
off one null set. Their limit is
\((1-\beta)X_\infty=\alpha X_\infty\).
The two-point set \(\{0,\alpha\}\) is closed, so
\(\alpha X_\infty\) belongs to it. Since \(\alpha>0\),
\(X_\infty\in\{0,1\}\).
Consequently \(P(X_\infty=1)=EX_\infty=EX_0\).

The hypotheses are realizable: take independent uniform variables
\(U_{n+1}\), independent of \(X_0\), and define
\[
           X_{n+1}=\beta X_n+\alpha\mathbf1_{\{U_{n+1}\leq X_n\}}.
\]
This remains in \([0,1]\) and has exactly the two required conditional
probabilities relative to its input filtration, including \(X_n=0,1\).
\end{proof}
\end{annexitem}

\begin{annexitem}{annex:p15.backward-martingale-convergence}{thm:p15-backward-martingale-convergence}{Finite reversed strings, measurable tail limits and mean convergence}
\begin{proof}
The tower identity proves by induction that
\(Y_n=E(Y_0\mid\mathcal G_n)\). Hence the conditional-family theorem
gives uniform integrability and \(E|Y_n|\leq E|Y_0|\).

Fix rational \(a<b\). For every \(n\), the finite process
\[
               Y_n,Y_{n-1},\ldots,Y_0
       \quad\hbox{with filtration}\quad
               \mathcal G_n,\mathcal G_{n-1},\ldots,\mathcal G_0
\]
is a forward martingale: the backward conditional identity is precisely
its one-step forward identity. If desired extend it constantly after
its last time to use the nonnegative infinite-time notation.
Let \(U_n(a,b)\) be the maximum number of completed upcrossings
of \([a,b]\) in this finite string.
The greedy upcrossing construction already proved in
\cref{thm:p15-upcrossing-inequality} realizes this maximum.
That inequality gives
\[
       (b-a)EU_n(a,b)
             \leq E(Y_0-a)^+-E(Y_n-a)^+
             \leq E(Y_0-a)^+<\infty.
\]
Appending a new observation at the beginning of a string preserves all
of its previous crossing pairs. Thus \(U_n(a,b)\) is nondecreasing in \(n\).
Monotone convergence shows
\(E\lim_n U_n(a,b)\leq E(Y_0-a)^+/(b-a)<\infty\);
the limiting crossing count is finite almost surely.

Suppose at an outcome that
\(\liminf_nY_n<\limsup_nY_n\), with extended limits allowed.
Choose rational \(a<b\) strictly between them.
There are arbitrarily late values at least \(b\) and at most \(a\).
For every \(r\), choose successive indices
\(k_1<k_2<\cdots<k_{2r}\) whose values alternate high, low,
starting high. Read the observations backward from \(k_{2r}\)
to zero. The selected low-high pairs now give at least \(r\)
upcrossings. Therefore \(U_n(a,b)\to\infty\), a contradiction.
The countable family of rational pairs consequently proves existence
of an extended limit almost surely.
Fatou's lemma and the first-moment bound give
\[
                        E\liminf_n|Y_n|\leq E|Y_0|<\infty.
\]
An infinite limiting magnitude thus has probability zero.
The limit is real and integrable.

There is one version measurable in the actual intersection of the
sigma-algebras. For each fixed \(k\), every \(Y_n\), \(n\geq k\),
is \(\mathcal G_k\)-measurable. The event of finite convergence and
the tail limsup are therefore measurable in every \(\mathcal G_k\).
Define \(Y_\infty\) to equal the finite limit on that event and zero
elsewhere. It is \(\mathcal G_\infty\)-measurable without a completeness
assumption. Uniform integrability and Vitali give \(Y_n\to Y_\infty\)
in \(L^1\).

If \(A\in\mathcal G_\infty\), then \(A\in\mathcal G_n\) for every \(n\).
The identity \(Y_n=E(Y_0\mid\mathcal G_n)\) gives
\(\int_A Y_n\,dP=\int_A Y_0\,dP\).
Convergence in \(L^1\) passes this equality to \(Y_\infty\).
Together with its measurability and integrability, it identifies
\(Y_\infty=E(Y_0\mid\mathcal G_\infty)\) almost surely.
\end{proof}
\end{annexitem}

\begin{annexitem}{annex:p15.levy-backward}{thm:p15-levy-backward}{Downward conditioning and the finite-coordinate extension}
\begin{proof}
Set \(Y_n=E(Z\mid\mathcal G_n)\). Integrability and adaptedness
hold by construction, and the tower identity gives
\[
                   E(Y_n\mid\mathcal G_{n+1})
                         =E(Z\mid\mathcal G_{n+1})=Y_{n+1}.
\]
This is a backward martingale. Its almost-sure and \(L^1\) limit is
\[
 E(Y_0\mid\mathcal G_\infty)
      =E(E(Z\mid\mathcal G_0)\mid\mathcal G_\infty)
      =E(Z\mid\mathcal G_\infty),
\]
where the last equality again uses the tower identity and
\(\mathcal G_\infty\subseteq\mathcal G_0\).

For a vector with finitely many integrable coordinates, apply the scalar
upward or downward theorem to each coordinate. The finite intersection
of their probability-one convergence events still has probability one.
For convergence in mean in, for example, the sum norm, the expectation
of the norm difference is the finite sum of the coordinate \(L^1\)
differences, which tends to zero. Other fixed finite-dimensional norms
are bounded by a constant times this norm, by the triangle inequality
applied to the coordinate unit vectors. This proves the stated vector
extension with all null-set and integrability requirements.
\end{proof}
\end{annexitem}
